% !TEX root = master.tex
% !TEX program = lualatex
% =========================
% Résumé complet — C01 (EPFL CS-202)
% =========================
\section{C01 — C Basics (résumé complet)}

\subsection{Objectifs du cours}
\begin{itemize}
  \item Apprendre à programmer \textbf{plus proche du système}, entre Java et l’assembleur.
  \item Savoir écrire des programmes C qui \textbf{manipulent la mémoire}, \textbf{utilisent des fichiers} et \textbf{les arguments de la ligne de commande}. 
  \item Ce cours vise des personnes qui savent déjà programmer (souvent Java) et insiste sur les différences/subtilités.
\end{itemize}

\subsection{Le langage C (positionnement général)}
\begin{itemize}
  \item C est un langage \textbf{typé}, \textbf{impératif}, \textbf{compilé}, \textbf{proche de la machine}.
  \item Différences fréquentes avec Java : pas de GC, moins de garde-fous (accès mémoire), beaucoup plus de responsabilité côté programmeur.
\end{itemize}

\subsection{Structure d’un programme C}
Structure typique :
\begin{itemize}
  \item \texttt{\#include} (bibliothèques)
  \item (éviter) variables globales
  \item prototypes
  \item \texttt{main}
\end{itemize}

\begin{lstlisting}[language=C]
#include <stdio.h>

int main(void) {
  printf("Hello\n");
  return 0;
}
\end{lstlisting}

\subsection{Variables et initialisation}
\begin{itemize}
  \item En C, une variable est \textbf{modifiable par défaut}.
  \item \textbf{Très important : pas d’initialisation automatique} (contrairement à Java) → \textbf{toujours initialiser}.
  \item Initialisation lors de la déclaration : \texttt{type ident = valeur;}
\end{itemize}

\begin{lstlisting}[language=C]
int val = 2;
double pi = 3.1415;
char c = 'a';
\end{lstlisting}

Valeurs littérales/échappements utiles (caractères) : \texttt{\textbackslash 0}, \texttt{\textbackslash n}, \texttt{\textbackslash \textbackslash}, \texttt{\textbackslash '} 

\subsection{Const (donnée non modifiable via ce nom)}
\begin{itemize}
  \item \texttt{const} signifie « \textit{read-only via ce nom} », pas « immuable pour toujours » (un pointeur peut modifier la zone mémoire). 
\end{itemize}

\begin{lstlisting}[language=C]
int const couple = 2;
double const g = 9.81;
\end{lstlisting}

\subsection{Types et modificateurs (portabilité)}
Types de base : \texttt{int}, \texttt{double}, \texttt{char}.  
Types « avancés » via modificateurs : \texttt{short}, \texttt{long}, \texttt{unsigned}, \texttt{long long}.

\begin{itemize}
  \item Depuis C99 : \texttt{bool} (\texttt{stdbool.h}), entiers \texttt{int8\_t}… (\texttt{stdint.h}), \texttt{long long}, \texttt{double complex} (\texttt{complex.h}). 
  \item La norme ne fixe pas les tailles exactes : seulement des inégalités comme
  \[
    \texttt{char} \le \texttt{short} \le \texttt{int} \le \texttt{long}, \quad \texttt{double} \le \texttt{long double}
  \]
  → attention à la portabilité. 
\end{itemize}

\subsection{Affectation : différence clé avec Java}
\begin{itemize}
  \item En C, \texttt{=} copie la \textbf{valeur} (et la sémantique diffère fortement de Java). 
\end{itemize}

\subsection{Opérateurs et expressions (rappels essentiels)}
\begin{itemize}
  \item Arithmétiques : \texttt{+ - * / \%}, incr./décr. \texttt{++ --}, composés \texttt{+=} etc.
  \item Comparaison : \texttt{== != < > <= >=}
  \item Logiques : \texttt{\&\& || !} (short-circuit)
  \item Bit-à-bit : \texttt{\& | \^{} \~{} << >>}
\end{itemize}

\textbf{Piège classique : division entière}
\begin{lstlisting}[language=C]
int a = 5 / 2;    // 2
double b = 5.0/2; // 2.5
\end{lstlisting}

\subsection{Structures de contrôle}
\subsubsection{Blocs et portée}
\begin{lstlisting}[language=C]
{
  int x = 3;   // visible seulement ici
}
\end{lstlisting}

\subsubsection{if / else}
\begin{lstlisting}[language=C]
if (cond) {
  ...
} else {
  ...
}
\end{lstlisting}

\subsubsection{switch (souvent avec enum)}
\begin{lstlisting}[language=C]
switch (x) {
  case 1: ...; break;
  default: ...; break;
}
\end{lstlisting}
Sans \texttt{break} : \textit{fall-through}.

\subsubsection{Boucles}
\begin{lstlisting}[language=C]
while (cond) { ... }

do { ... } while (cond);

for (int i = 0; i < n; ++i) { ... }
\end{lstlisting}

\texttt{break} : sort de la boucle, \texttt{continue} : saute à l’itération suivante.

\subsection{Fonctions : définition, prototype, passage par valeur/référence}
\begin{itemize}
  \item Déclaration (prototype) avant usage.
  \item Passage \textbf{par valeur} par défaut.
  \item Pour modifier une variable de l’appelant : passer un \textbf{pointeur} (adresse) — notion de « par référence » (valeur de pointeur en fait).
\end{itemize}

\begin{lstlisting}[language=C]
int f(int x) {         // passage par valeur
  x = x + 1;
  return x;
}

void g(int* px) {      // passage par "référence" via pointeur
  *px = *px + 1;
}
\end{lstlisting}

\subsection{Tableaux (taille fixe en C)}
\begin{itemize}
  \item En C, un tableau est \textbf{homogène} et de \textbf{taille fixe}. 
  \item Indices : \texttt{0} à \texttt{taille-1}; \textbf{pas de contrôle de débordement} → buffer overflow possible.
  \item Un tableau \textbf{ne connaît pas sa taille} : il faut la gérer séparément. 
\end{itemize}

\begin{lstlisting}[language=C]
#define N 5
int age[N];

for (size_t i = 0; i < N; ++i) {
  printf("Age %zu ? ", i+1);
  scanf("%d", &age[i]);
}
\end{lstlisting}

\subsubsection{Initialisation}
\begin{lstlisting}[language=C]
int t[3] = { 1, 2, 3 };
\end{lstlisting}

\subsubsection{Tableaux multidimensionnels}
Déclaration \texttt{type tab[a][b];} et accès \texttt{tab[i][j]}. 

\begin{lstlisting}[language=C]
double rotation[2][2];
rotation[0][1] = 3.14;
\end{lstlisting}

\subsubsection{Passage d’un tableau à une fonction}
\begin{itemize}
  \item Un paramètre \texttt{type tab[]} est en réalité traité comme \texttt{type* tab}.
  \item Donc l’appel peut modifier les éléments (effets de bord) : ajouter \texttt{const} si nécessaire.
  \item Conseil : toujours passer la taille en plus. 
\end{itemize}

\begin{lstlisting}[language=C]
void fill(int tab[], size_t n) {
  for (size_t i = 0; i < n; ++i) tab[i] = (int)(i+1);
}

void print(int const tab[], size_t n) { // pas de modif
  for (size_t i = 0; i < n; ++i) printf("%d ", tab[i]);
}
\end{lstlisting}

\subsection{Types énumérés (enum)}
\begin{itemize}
  \item Sert à donner des noms à des valeurs entières (couleurs, cantons, etc.). 
  \item Par convention, le premier vaut 0, puis 1, 2, … 
  \item Pratique avec \texttt{switch} et indexation de tableaux.
\end{itemize}

\begin{lstlisting}[language=C]
enum CantonRomand { Vaud, Valais, Geneve, Neuchatel, Fribourg, Jura };

enum CantonRomand c = Vaud;
switch (c) {
  case Vaud:   /* ... */ break;
  case Valais: /* ... */ break;
  default:     /* ... */ break;
}
\end{lstlisting}

\subsection{typedef (alias de types)}
\begin{itemize}
  \item Permet de renommer un type (souvent pour lisibilité, maintenance, concepts). 
  \item Très utile pour tableaux et signatures de fonctions. 
\end{itemize}

\begin{lstlisting}[language=C]
typedef unsigned long int Compteur;
typedef double Matrice3[3][3];

#define N 3
typedef double Vecteur[N];
typedef Vecteur Matrice[N];
\end{lstlisting}

\subsection{Structures (struct)}
\begin{itemize}
  \item Pour regrouper des champs hétérogènes (contrairement aux tableaux homogènes). 
  \item Déclaration : \texttt{struct Nom \{ ... \};} puis \texttt{struct Nom var;}.
  \item Accès aux champs : \texttt{.} ; si pointeur : \texttt{->}.
\end{itemize}

\begin{lstlisting}[language=C]
#define TAILLE_MAX_NOM 64

struct Personne {
  char nom[TAILLE_MAX_NOM];
  double taille;
  int age;
  char sexe;
};

struct Personne p;
p.age = 20;
\end{lstlisting}

\subsubsection{typedef + struct (style pratique)}
\begin{lstlisting}[language=C]
typedef struct {
  char nom[TAILLE_MAX_NOM];
  double taille;
  int age;
  char sexe;
} Personne;

Personne untel;
\end{lstlisting}
(Plus besoin du mot \texttt{struct}.) 

\subsubsection{Initialisation de struct}
\begin{lstlisting}[language=C]
Personne a = { "Dupontel", 1.75, 20, 'M' };
\end{lstlisting}
Et depuis C99 : initialisation par champs (\textit{designated initializers}).

\subsubsection{Affectation de structures}
Une struct peut être affectée globalement (\texttt{p2 = p1;}) : copie champ par champ.

\subsubsection{Exemple complet (struct + I/O)}
Inspiré des slides : affichage, modification via pointeur, création via saisie. 

\begin{lstlisting}[language=C]
#include <stdio.h>
#include <string.h>

typedef struct {
  double taille;
  int age;
  char sexe;
} Personne;

void affiche_personne(Personne p) {
  printf("taille: %f\n", p.taille);
  printf("age   : %d\n", p.age);
  printf("sexe  : %c\n", p.sexe);
}

void anniversaire(Personne* p) {
  ++(p->age);
}

Personne naissance(void) {
  Personne p;
  memset(&p, 0, sizeof(p));
  puts("Saisie d'une nouvelle personne");
  printf("Entrez sa taille : ");
  scanf("%lf", &p.taille);
  printf("Entrez son age   : ");
  scanf("%d", &p.age);
  printf("Homme [M] ou Femme [F] : ");
  scanf(" %c", &p.sexe); // espace pour ignorer blancs
  return p;
}
\end{lstlisting}

\subsection{Autres types / remarques importantes}
\begin{itemize}
  \item En C, il n’y a pas de type \texttt{string} natif : une chaîne est un \texttt{char*} (ou \texttt{char[]} terminé par \texttt{'\textbackslash 0'}).
  \item Recommandation sécurité : toujours borner les tailles lors des lectures/copies (préférer \texttt{strncpy} à \texttt{strcpy}) → éviter \textit{buffer overflow}.
\end{itemize}

\subsection{Entrée/Sortie standard (stdio.h)}
Flots standards :
\begin{itemize}
  \item \texttt{stdin} (clavier/entrée), \texttt{stdout} (sortie), \texttt{stderr} (erreurs).
  \item \texttt{stderr} est séparé et non bufferisé → utile pour messages d’erreur. 
\end{itemize}

\subsubsection{printf : formats et options}
\begin{itemize}
  \item \texttt{printf} retourne le nombre de caractères écrits (ou négatif si échec). 
  \item Options entre \% et la conversion : \texttt{- + ' ' \# 0}, largeur, \texttt{.precision}, modificateurs \texttt{h l L}.
  \item Conversions utiles : \texttt{\%d \%u \%x \%o \%c \%s \%f \%e \%g \%p \%\%}.
\end{itemize}

\begin{lstlisting}[language=C]
double x = 10.4276;
printf(">%7.2f%%<\n", x);   // largeur 7, precision 2, affiche un %
\end{lstlisting}

\subsubsection{Buffering et fflush}
\begin{itemize}
  \item \texttt{printf} écrit souvent dans un buffer : un message peut ne pas être affiché avant un crash.
  \item On peut forcer l’affichage avec \texttt{fflush(stdout)}. 
  \item \texttt{fflush(stdin)} n’a pas de sens (aucun effet utile).
\end{itemize}

\begin{lstlisting}[language=C]
printf("Entrez un nombre: ");
fflush(stdout);
\end{lstlisting}

\subsubsection{scanf : principe + pièges}
\begin{itemize}
  \item \texttt{scanf("FORMAT", ptr1, ptr2, ...)} stocke dans la mémoire pointée (souvent \texttt{\&var}). 
  \item Les blancs (\texttt{isspace}) séparent les valeurs (et peuvent être absorbés). 
  \item Pour lire une ligne avec des espaces : préférer \texttt{fgets} (contrôle de taille, \texttt{\textbackslash 0} garanti).
  \item \textbf{Toujours} contrôler la valeur de retour de \texttt{scanf} (nb de champs convertis, ou \texttt{EOF}). 
\end{itemize}

\paragraph{Formats de lecture (spécifiques scanf).}
\begin{itemize}
  \item \texttt{\%lf} pour lire un \texttt{double}.
  \item \texttt{\%ns} (ex. \texttt{\%19s}) pour borner la taille lue.
  \item \texttt{\%[A-Z]} / \texttt{\%[^\textbackslash n]} : lecture par ensembles de caractères. 
\end{itemize}

\paragraph{Piège : boucle infinie si saisie invalide.}
\begin{itemize}
  \item Exemple : si on attend un int et que l’utilisateur tape \texttt{a}, on peut boucler infiniment.
  \item Solution : vérifier le retour de \texttt{scanf}, puis vider le buffer d’entrée (avec \texttt{getc} jusqu’au \texttt{\textbackslash n}).
\end{itemize}

\begin{lstlisting}[language=C]
#include <stdio.h>

int main(void) {
  int i = 0, j = 0;
  do {
    printf("Entrez un nombre entre 1 et 10 : ");
    fflush(stdout);
    j = scanf("%d", &i);

    if (j != 1) {
      printf("Je vous ai demandé un nombre, pas du charabia !\n");
      while (!feof(stdin) && !ferror(stdin) && getc(stdin) != '\n') { }
    }
  } while (!feof(stdin) && !ferror(stdin) && ((j != 1) || (i < 1) || (i > 10)));

  return 0;
}
\end{lstlisting}

\subsubsection{fgets (lecture de ligne)}
\begin{itemize}
  \item \texttt{fgets(buf, size, stdin)} lit au plus \texttt{size-1} caractères et garantit le \texttt{'\textbackslash 0'}.
  \item Si un \texttt{'\textbackslash n'} est lu, il est inclus dans la chaîne → on peut le supprimer ensuite.
\end{itemize}

\subsubsection{sprintf / sscanf (flots dans des chaînes)}
\begin{itemize}
  \item Pour écrire dans une chaîne (au lieu de stdout) : \texttt{sprintf}.
  \item \texttt{sscanf} existe aussi (conversion depuis une chaîne).
\end{itemize}

\begin{lstlisting}[language=C]
char adresse[129];
int code = 1000;
char ville[] = "Lausanne";
sprintf(adresse, "%5d %.122s", code, ville);
\end{lstlisting}

\subsection{Fichiers (FILE*)}
Résumé de la section fichiers. 

\subsubsection{Ouverture et modes}
\begin{itemize}
  \item Type : \texttt{FILE*}.
  \item \texttt{FILE* fopen(const char* nom, const char* mode)}.
  \item Modes : \texttt{"r"} lecture, \texttt{"w"} écriture (écrase), \texttt{"a"} append; \texttt{+} pour lecture+écriture; \texttt{b} pour binaire. 
\end{itemize}

\begin{lstlisting}[language=C]
FILE* f1 = fopen("in.txt", "r");
FILE* f2 = fopen("out.txt", "w");
FILE* f3 = fopen("data.bin", "a+b");
\end{lstlisting}

\subsubsection{Gestion d’erreur à l’ouverture}
\begin{itemize}
  \item Si \texttt{fopen} échoue → retourne \texttt{NULL}.
  \item Pour plus d’infos : \texttt{perror}, \texttt{strerror(errno)}.
\end{itemize}

\begin{lstlisting}[language=C]
#include <stdio.h>

FILE* f = fopen("inexistant.txt", "r");
if (f == NULL) {
  perror("fopen");
  return 1;
}
\end{lstlisting}

\subsubsection{Fichiers texte : fprintf / fscanf}
Principe : mêmes formats que \texttt{printf/scanf} mais sur un \texttt{FILE*}.

\begin{lstlisting}[language=C]
FILE* entree = fopen("in.txt", "r");
FILE* sortie = fopen("out.txt", "a");
int i = 0;

if (entree && sortie) {
  fscanf(entree, "%d", &i);
  fprintf(sortie, "%d\n", i);
}
\end{lstlisting}

\subsubsection{Fin de fichier et erreurs}
\begin{itemize}
  \item \texttt{feof(FILE*)} : fin de fichier atteinte (non nul).
  \item \texttt{ferror(FILE*)} : erreur sur le fichier (non nul).
\end{itemize}

\begin{lstlisting}[language=C]
while (!feof(entree) && !ferror(entree)) {
  /* lecture */
}
\end{lstlisting}

\subsubsection{Fermeture}
\begin{itemize}
  \item \texttt{fclose(FILE*)} ferme un flot. \textbf{Ne pas oublier} surtout après écriture. 
\end{itemize}

\subsubsection{Exemple : lecture fichier texte mot par mot}
Exemple directement inspiré des slides (borne \%10s). 
\begin{lstlisting}[language=C]
#include <stdio.h>

int main(void) {
  FILE* entree = fopen("test", "r");
  if (entree == NULL) {
    fprintf(stderr, "Erreur : impossible de lire le fichier %s\n", "test");
    return 1;
  }

  #define TAILLE_MAX 10
  char mot[TAILLE_MAX+1] = "";

  while (fscanf(entree, "%10s", mot) == 1) {
    printf("lu : ->%-*s<-\n", TAILLE_MAX, mot);
  }

  fclose(entree);
  return 0;
}
\end{lstlisting}

\subsubsection{Écriture dans un fichier avec fgets/fputs}
Exemple des slides : demander un nom de fichier, lire une phrase, écrire.
\begin{lstlisting}[language=C]
#include <stdio.h>
#include <string.h>

int main(void) {
  char nom_fichier[FILENAME_MAX+1] = "";
  FILE* sortie = NULL;

  size_t len = 0;
  do {
    printf("Dans quel fichier voulez-vous écrire ?\n");
    fgets(nom_fichier, FILENAME_MAX+1, stdin);
    len = strlen(nom_fichier);
    if (len != 0 && nom_fichier[len-1] == '\n') {
      nom_fichier[--len] = '\0';
    }
  } while (len == 0 && !feof(stdin) && !ferror(stdin));

  if (nom_fichier[0] != '\0') {
    sortie = fopen(nom_fichier, "a");
    if (sortie == NULL) {
      fprintf(stderr, "Erreur : impossible d'écrire dans %s\n", nom_fichier);
    } else {
      #define MAX_PHRASE 513
      char phrase[MAX_PHRASE] = "";
      printf("Entrez une phrase :\n");
      fgets(phrase, MAX_PHRASE, stdin);
      fputs(phrase, sortie);
      fclose(sortie);
    }
  }
  return 0;
}
\end{lstlisting}

\subsubsection{Flots standards vs fichiers}
\begin{itemize}
  \item \texttt{printf} = \texttt{fprintf(stdout, ...)}
  \item \texttt{scanf} = \texttt{fscanf(stdin, ...)}
  \item Messages d’erreur : \texttt{fprintf(stderr, ...)} 
\end{itemize}

\subsubsection{Fichiers binaires : fwrite / fread}
\begin{itemize}
  \item \texttt{fwrite(ptr, taille\_el, nb\_el, fichier)} écrit des données brutes; retourne nb d’éléments écrits.
  \item \texttt{fread(ptr, taille\_el, nb\_el, fichier)} lit des données brutes; retourne nb d’éléments lus.
\end{itemize}

\paragraph{Écriture binaire (exemple).} 
\begin{lstlisting}[language=C]
#include <stdio.h>
#define TAILLE 12

int main(void) {
  int tab[TAILLE];
  for (size_t i = 0; i < TAILLE; ++i) tab[i] = (int)(i*i - 6*i);

  FILE* sortie = fopen("test-fwrite.bin", "wb");
  if (sortie != NULL) {
    size_t nb_ok = fwrite(tab, sizeof(int), TAILLE, sortie);
    fclose(sortie);
    if (nb_ok != TAILLE) {
      fprintf(stderr, "Unable to write %zu elements, wrote only %zu\n",
              (size_t)TAILLE, nb_ok);
    }
  }
  return 0;
}
\end{lstlisting}

\paragraph{Lecture binaire (exemple).} 
\begin{lstlisting}[language=C]
#include <stdio.h>
#define TAILLE 12

int main(void) {
  FILE* entree = fopen("test-fwrite.bin", "rb");
  if (entree != NULL) {
    int tab[TAILLE];
    size_t nb_ok = fread(tab, sizeof(int), TAILLE, entree);
    fclose(entree);

    if (nb_ok != TAILLE) {
      fprintf(stderr, "Unable to read %zu elements, read only %zu\n",
              (size_t)TAILLE, nb_ok);
    } else {
      for (size_t i = 0; i < TAILLE; ++i)
        printf("tab[%2zu] = % 3d\n", i, tab[i]);
    }
  }
  return 0;
}
\end{lstlisting}

\subsubsection{Manipulations diverses (fseek/ftell/rewind/clearerr)}
\begin{itemize}
  \item \texttt{fseek(FILE*, offset, SEEK\_SET|SEEK\_CUR|SEEK\_END)} repositionne la tête de lecture.
  \item \texttt{ftell(FILE*)} donne la position courante.
  \item \texttt{rewind(FILE*)} revient au début.
  \item \texttt{clearerr(FILE*)} remet à zéro l’état d’erreur.
\end{itemize}

\subsection{Règles d’or (sécurité / robustesse)}
\begin{itemize}
  \item Toujours borner les lectures (\texttt{\%ns}, \texttt{fgets(size)}), et contrôler les retours (\texttt{scanf}, \texttt{fscanf}, \texttt{fread}, \texttt{fwrite}). 
  \item Ne jamais oublier \texttt{fclose}.
  \item Préférer \texttt{stderr} pour erreurs.
  \item Attention aux débordements (buffer overflow) : tailles, copies (\texttt{strncpy} vs \texttt{strcpy}).
\end{itemize}
